% =============================================================================
%  fault_location_id - Appendix B
%
%  WP1.6 deliverable.  Standalone derivation of the corrected
%  Cramér-Rao bound for the HIF-TF locator under the dual-channel
%  AWGN noise model + ratio observation.  Closes risk R1 (Gaussian-
%  on-H FIM invalid in the HIF regime) and risk R9 (Geary-Hinkley
%  validity reporting).  Builds with `pdflatex AppendixB...tex` x 2.
% =============================================================================

\documentclass[11pt,a4paper]{article}

\usepackage[T1]{fontenc}
\usepackage[utf8]{inputenc}
\usepackage[margin=22mm]{geometry}
\usepackage{amsmath,amssymb,amsfonts,bm}
\usepackage{booktabs}
\usepackage{siunitx}
\sisetup{detect-all,per-mode=symbol}
\usepackage{enumitem}
\usepackage{titlesec}
\usepackage{xcolor}
\definecolor{sambpblue}{HTML}{1F3D7A}
\titleformat{\section}
  {\color{sambpblue}\Large\bfseries}{Appendix~B.\thesection}{0.6em}{}
\titleformat{\subsection}
  {\color{sambpblue}\large\bfseries}{B.\thesection.\arabic{subsection}}{0.6em}{}
\usepackage[hidelinks,colorlinks=true,linkcolor=sambpblue,urlcolor=sambpblue]{hyperref}

\title{Appendix B.\\[2pt]
       Corrected Cram\'er--Rao Lower Bound\\
       under the Dual-Channel AWGN Noise Model\\[6pt]
       \large for the SAMBPS DTaaS HIF-TF Locator (\texttt{fault\_location\_id})}
\author{Arjundas K., K. Shanthi Swarup\\[1pt]
        \small Power Systems Computational Lab, IIT Madras}
\date{Issued under WP1.6; v3 Execution Plan deliverable D-B; closes R1, R9}

\begin{document}
\maketitle

\noindent
The IEEE Access v1 manuscript reports a Cram\'er--Rao analysis under
a circular complex Gaussian noise model placed directly on the
single-bin admittance $H_{\mathrm{meas}} = I_{\mathrm{bin}}/V_{\mathrm{bin}}$.
That model is \emph{not} self-consistent with \S\,V-B of the same
manuscript, which specifies dual-channel AWGN on the time-domain
voltage and current channels.  $H_{\mathrm{meas}}$ is the ratio of
two independent complex Gaussians whose density is the
Marsaglia--Nadarajah--Pog\'any form, not Gaussian.  This appendix
derives two corrected bounds:

\begin{enumerate}
  \item the \textbf{proper-complex-Gaussian-ratio CRLB} after
        Kuruo\u{g}lu (2018), which uses the high-SNR$_V$ Gaussian
        approximation of the ratio density and includes the
        Geary--Hinkley validity test;
  \item the \textbf{joint dual-channel FIM} in $(V(t), I(t))$ space
        after Nehorai \& Hawkes (2000), projected onto $(\alpha, R_x)$.
\end{enumerate}

The two are cross-checked at high SNR$_V$ and used to overlay the
empirical RMS error from the WP1.5 Monte-Carlo on the
$2\times 4$ panel \texttt{outputs/phase1\_crlb\_overlay/}.

% =============================================================================
\section{Ratio-density form (Marsaglia--Nadarajah--Pog\'any)}
% =============================================================================

For independent complex Gaussian random variables
$X = \mu_X + n_X$ and $Y = \mu_Y + n_Y$ with circular complex
variances $\sigma_X^2$ and $\sigma_Y^2$ per real/imag component,
the ratio $Z = X/Y$ has a non-Gaussian density
\cite{Marsaglia1965Ratio,NadarajahPogany2018} that can be written
in closed form via Bessel functions but is unwieldy.  In the
high-SNR$_Y$ regime, defined by the Geary--Hinkley criterion
\begin{equation}
   \frac{|\mu_Y|}{\sigma_Y} > \tau, \quad \tau \approx 4,
   \label{eq:gh}
\end{equation}
a first-order Taylor expansion in $n_Y/\mu_Y$ gives
\begin{equation}
   Z \;\approx\; \frac{\mu_X}{\mu_Y}
                 \;+\; \frac{n_X}{\mu_Y}
                 \;-\; \frac{\mu_X}{\mu_Y^2}\,n_Y,
   \label{eq:taylor}
\end{equation}
which is itself a complex Gaussian with mean $\mu_X/\mu_Y$ and
\emph{signal-dependent} per-component variance
\begin{equation}
   \sigma_Z^2 \;=\; \frac{\sigma_X^2 + |Z|^2 \sigma_Y^2}{|\mu_Y|^2}.
   \label{eq:sigmaZ}
\end{equation}
The $|Z|^2 \sigma_Y^2$ term is the ``ratio-shot'' contribution that
the v1 manuscript drops by treating $\sigma_H^2$ as constant.

% =============================================================================
\section{Proper-ratio FIM construction}
% =============================================================================

Substituting $X = I_{\mathrm{bin}}$, $Y = V_{\mathrm{bin}}$,
$Z = H_{\mathrm{meas}}$ in (\ref{eq:sigmaZ}):
\begin{equation}
   \sigma_H^2 \;=\;
       \frac{\sigma_I^2 + |H|^2\,\sigma_V^2}{|V_{\mathrm{phase}}|^2},
   \label{eq:sigmaH}
\end{equation}
where $|V_{\mathrm{phase}}|$ is the phasor magnitude of the
ideal source at $\omega_0$ and $\sigma_X^2 = 2\sigma_{x_t}^2 / N_s$
is the per-real/imag variance of the single-bin DFT phasor of a
length-$N_s$ time-domain channel with white Gaussian noise of
variance $\sigma_{x_t}^2$.

For complex observations $h(\theta) = H(\theta)$ contaminated by
circular complex Gaussian noise of per-component variance $\sigma_H^2$,
the Fisher information matrix on $\theta = (\alpha, R_x)$ is
\cite{Kay1993Vol1}
\begin{equation}
   F^{\mathrm{proper}}_{kl}(\theta) \;=\; \frac{1}{\sigma_H^2}\,
         \mathrm{Re}\!\left[
           \left(\partial_{\theta_k} H\right)^{\!*}
           \,\partial_{\theta_l} H
         \right].
   \label{eq:Fproper}
\end{equation}
The CRLB on the variance of any unbiased estimator of $\theta_k$ is
$\bigl(F^{\mathrm{proper}}\bigr)^{-1}_{kk}$.  Substituting
(\ref{eq:sigmaH}):
\begin{equation}
   F^{\mathrm{proper}}_{kl}
       \;=\;
       \frac{|V_{\mathrm{phase}}|^2}
            {\sigma_I^2 + |H|^2 \sigma_V^2}\;
       \mathrm{Re}\!\left[
         \left(\partial_{\theta_k} H\right)^{\!*}
         \,\partial_{\theta_l} H
       \right].
   \label{eq:FproperFinal}
\end{equation}
The partials $\partial H/\partial \alpha$ and $\partial H/\partial R_x$
are computed by the central-finite-difference fallback in
\texttt{inverse\_estimation/faultloc\_crlb\_proper.py} and will be
swapped for the closed-form Phase-2 forms (WP2.2) when those land.

\paragraph{Geary--Hinkley validity.}  $\sigma_H^2$ in
(\ref{eq:sigmaH}) is the variance of a Gaussian \emph{approximation}
of the Marsaglia density.  Equation (\ref{eq:sigmaH}) is a faithful
approximation only when the V channel is in the Geary--Hinkley regime
(\ref{eq:gh}).  Outside that regime the true ratio density carries
heavier tails than the Gaussian, and (\ref{eq:FproperFinal})
\emph{overestimates} the FIM (i.e.\ underestimates the true CRLB).
The library reports a per-cell boolean \texttt{geary\_hinkley\_valid}
flag.  Under the project's $N_s = 200$ single-bin averaging the GH
threshold $|V_{\mathrm{phase}}|/\sigma_V > 4$ corresponds to
time-domain $\mathrm{SNR}_V \gtrsim -8$\,dB; the entire 720-case
grid (which starts at 20\,dB) is comfortably in the regime.

% =============================================================================
\section{Dual-channel FIM (Nehorai--Hawkes)}
% =============================================================================

The dual-channel FIM treats the raw waveform vector
$\mathbf{y} = [v(t_0),\ldots,v(t_{N_s-1}),\,i(t_0),\ldots,i(t_{N_s-1})]^T$
as the observation, with
$v(t_n) = V_{\mathrm{phase}}\cos(\omega_0 t_n) + n_v(n)$ and
$i(t_n) = \mathrm{Re}\!\left\{H(\theta) V_{\mathrm{phase}}
e^{j\omega_0 t_n}\right\} + n_i(n)$.  Under the project's
ideal-source assumption $\partial v/\partial \theta \equiv 0$, so the
V channel contributes zero to the FIM; only the I channel carries
information about $\theta$.  Defining the Jacobian
$[J_i]_{n,k} = \partial i_{\mathrm{clean}}(t_n)/\partial \theta_k$,
\begin{equation}
   F^{\mathrm{dual}}_{kl}(\theta)
       \;=\; \frac{1}{\sigma_{i_t}^2}
             \sum_{n=0}^{N_s-1}
                 \frac{\partial i_{\mathrm{clean}}(t_n)}{\partial \theta_k}\,
                 \frac{\partial i_{\mathrm{clean}}(t_n)}{\partial \theta_l}.
   \label{eq:Fdual}
\end{equation}
A one-cycle window with $N_s$ samples evaluates the inner sum
analytically (cosine--cosine and sine--sine sums each contribute
$N_s/2$, cross terms cancel):
\begin{equation}
   F^{\mathrm{dual}}_{kl}
        \;=\;
        \frac{N_s}{2}\,\frac{V_{\mathrm{phase}}^2}{\sigma_{i_t}^2}\,
        \mathrm{Re}\!\left[
          \left(\partial_{\theta_k} H\right)^{\!*}
          \,\partial_{\theta_l} H
        \right].
   \label{eq:FdualFinal}
\end{equation}

% =============================================================================
\section{Cross-check: proper-ratio vs dual-channel}
% =============================================================================

The same Fisher kernel $\mathrm{Re}\!\left[(\partial_k H)^{*}
\,\partial_l H\right]$ appears in both (\ref{eq:FproperFinal}) and
(\ref{eq:FdualFinal}); the prefactors differ.  Substituting
$\sigma_I^2 = 2\sigma_{i_t}^2/N_s$ and $\sigma_V^2 = 2\sigma_{v_t}^2/N_s$:
\begin{equation}
   \frac{F^{\mathrm{proper}}}{F^{\mathrm{dual}}}
        \;=\; \frac{\sigma_I^2}{\sigma_I^2 + |H|^2 \sigma_V^2}.
   \label{eq:ratio}
\end{equation}
Two limiting cases:
\begin{itemize}[leftmargin=2em]
  \item $\sigma_V \to 0$ (V perfectly known): the two FIMs are
        \emph{exactly equal}.  This is the \textbf{WP1.6 brief
        acceptance regime} (``agree to within 5\,\% at
        $\mathrm{SNR}_I \geq 40$\,dB in the Geary--Hinkley regime'')
        verified by
        \texttt{tests/test\_crlb\_consistency.py::test\_crlb\_proper\_eq\_dual\_when\_V\_noiseless}.
  \item $\sigma_V = \sigma_I/|H|$ (V and I shot noise balanced):
        $F^{\mathrm{proper}} = F^{\mathrm{dual}}/2$, i.e.\ a factor
        of two loss in information by accessing only the ratio
        instead of the raw waveforms.  In RMSE terms,
        $\mathrm{RMSE}^{\mathrm{proper}} = \sqrt{2}\,
        \mathrm{RMSE}^{\mathrm{dual}}$.  Verified by
        \texttt{test\_crlb\_factor\_sqrt2\_when\_snrV\_eq\_snrI}.
\end{itemize}

\paragraph{Direction of the inequality.}  Note $F^{\mathrm{proper}}
\leq F^{\mathrm{dual}}$ always, with equality iff $\sigma_V = 0$.
The proper-ratio CRLB is therefore $\geq$ the dual-channel CRLB
(it represents \emph{less} information about $\theta$ because the
H-ratio observation discards the absolute magnitudes that the raw V
and I waveforms preserve).  This is opposite to the v1
``Gaussian-on-H'' linearisation, which artificially \emph{tightens}
the bound by ignoring the $|H|^2 \sigma_V^2$ term in
(\ref{eq:sigmaH}).

% =============================================================================
\section{Headline empirical-vs-CRLB result}
% =============================================================================

The WP1.5 Monte-Carlo (100 trials per cell on the 720-case grid)
provides empirical per-cell $\mathrm{RMSE}(\hat\alpha)$ values that
the bounds (\ref{eq:FproperFinal}) and (\ref{eq:FdualFinal}) are
plotted against in the $2\times 4$ overlay panel
\texttt{outputs/phase1\_crlb\_overlay/}.  The four headline findings
this Appendix supports:

\begin{enumerate}[leftmargin=2em]
  \item The proper-ratio bound is the \emph{correct} bound on any
        unbiased estimator that operates on $H_{\mathrm{meas}}$.
        It supersedes the v1 ``Gaussian-on-H'' linearisation,
        which is invalid outside the asymptotic
        $|I| \gg \sigma_I$ regime.
  \item The dual-channel bound is the ground truth in $(V, I)$
        waveform space; it is \emph{tighter} (lower) than the
        proper-ratio bound by the factor in (\ref{eq:ratio}).
  \item The two bounds agree to within 5\,\% (in fact exactly) at
        $\sigma_V \to 0$, certified by
        \texttt{test\_crlb\_proper\_eq\_dual\_when\_V\_noiseless}.
  \item The Geary--Hinkley validity flag holds across the full
        720-case grid (time-domain $\mathrm{SNR}_V$ down to
        $-8$\,dB), so (\ref{eq:sigmaH}) is a faithful
        approximation throughout.  Outside this regime the proper-
        ratio FIM \emph{overestimates} the true Marsaglia FIM.
\end{enumerate}

\bigskip\hrule\bigskip
\noindent\textit{End of Appendix~B.}

% Bibliography stub (manuscript_v2.tex's references.bib carries the
% canonical entries; this Appendix relies on the parent .bib).
\bibliographystyle{IEEEtran}
\bibliography{references}

\end{document}
